% The PNAS significance statement as a standalone page, for the submission
% portal, which takes it as its own field rather than as part of the manuscript.
%
% The text here is the SINGLE SOURCE for that field, and it is kept identical
% to the \section*{Significance Statement} block of main.tex.  If one is edited
% the other has to follow; `make check-significance` compares them word for
% word and fails if they have drifted.
%
% PNAS: "Explain the significance of the research at a level understandable to
% an undergraduate-educated scientist outside their field of specialty.
% Include no more than 120 words."  This is 118.
\documentclass[11pt]{article}
\usepackage[left=1in,right=1in,top=3cm,bottom=0.35in,%
includeheadfoot,headheight=99pt,headsep=0.2in]{geometry}
\usepackage{times,graphicx,cite,fancyhdr,amsmath,amssymb,captdef}
\usepackage{marvosym}
\renewcommand{\headrulewidth}{0pt}
%----------------  page style -------------------
\pagestyle{fancy}
\fancyhead{}    %clear the head
\fancyhead[L]{
  \begin{minipage}[t]{1.8in}
    \hspace{0.6in}
    \includegraphics[height=0.6in]{anl.png}
  \end{minipage}
}
\fancyhead[C]{}
\fancyhead[R]{%
  \begin{minipage}[t]{1.7in}
    \begin{flushleft}
      {\scriptsize
        \vspace{-1.1in}
        {\bf Qingteng Zhang, Ph.D., Physicist}\\
        \vspace{0.08in}
        {\bf X-ray Science Division}\\
        Argonne National Laboratory \\
        9700 South Cass Avenue\\
        Argonne, IL 60439 \\
        \vspace{0.08in}
        \Telefon: 630-252-6366 \\
        }
    \end{flushleft}
  \end{minipage}
}
\fancyfoot{}    %clear the foot
\fancyfoot[C]{
  \footnotesize{A U.S. Department of Energy laboratory managed by UChicago
    Argonne, LLC}
}
% ==============================================================================
\setlength{\parskip}{2ex}
\setlength{\parindent}{0pt}

\usepackage[hidelinks]{hyperref}
\setlength{\parskip}{0.6em}
\setlength{\parindent}{0pt}
\raggedbottom

\begin{document}

{\large\bfseries Significance Statement}

\vspace{0.6em}

{\footnotesize\itshape
Temperature, Not Structure, Sets Dynamic Arrest in an Elastin-like Protein
Polymer \\
Ozgulbas, Chu, Lavens, Marks, Du, Zuo, Irvin, Ma, Seifert, Lee, Vescovi,
Foster, Mateevitsi, Jensen, Akins, Ramanathan, Babnigg and Zhang \\
Physical Sciences: Applied Physical Sciences; secondary, Biophysics and
Computational Biology}

\vspace{1.2em}

Whether a soft material has solidified is usually inferred from its structure:
a scattering pattern that has grown enough is read as a gel. We show that this
inference can fail. In a designed protein polymer, samples held within a degree
of one another develop nearly the same structure, yet only the warmer ones stop
moving. Seeing this needed X-ray speckle framed at fifty thousand images per
second, because the motion reporting on arrest is millisecond-scale, below what
conventional detectors reach. Arrest spreads region by region, largest scales
first, and one rescaling of waiting time collapses every sample onto a single
pathway whose temperature scale is a fifth of a degree, too sharp for single
contacts to set.

\vspace{1.4em}

{\footnotesize 118 words. Limit 120.}

\end{document}
